\documentclass[ecta,nameyear,draft]{econsocart}
\usepackage{amsmath,amsfonts,amssymb,amsthm,booktabs,array,longtable}
\usepackage[T1]{fontenc}
\usepackage{mathpazo,microtype,needspace,xurl}
\usepackage{xr-hyper}
\RequirePackage[colorlinks,citecolor=blue,linkcolor=blue,urlcolor=blue]{hyperref}
\startlocaldefs
\newtheorem{theorem}{Theorem}
\newtheorem{proposition}[theorem]{Proposition}
\newtheorem{lemma}[theorem]{Lemma}
\newtheorem{corollary}[theorem]{Corollary}
\newcommand{\E}{\mathbb E}
\newcommand{\R}{\mathbb R}
\endlocaldefs
\setlength{\emergencystretch}{3em}
\makeatletter\def\copyright@text{Prepared for referee review}\makeatother

\setcounter{secnumdepth}{2}

% Stable current-document identity for clean typesetting reproduction.
\ifdefined\pdfinfoomitdate\pdfinfoomitdate=1\fi
\ifdefined\pdfsuppressptexinfo\pdfsuppressptexinfo=15\fi
\ifdefined\pdftrailerid\pdftrailerid{}\fi

\externaldocument[][nocite]{build/supp}[supp.pdf]
\externaldocument[][nocite]{build/appendix}[appendix.pdf]
\begin{document}
\begin{frontmatter}
\title{Neural Bellman Operators}
\runtitle{Neural Bellman Operators}
\begin{aug}
\author[id=au1,addressref={add1}]{\fnms{Qian}~\snm{QI}\ead[label=e92ECTA]{qiqian@pku.edu.cn}}
\address[id=add1]{Peking University}
\end{aug}
\begin{abstract}
This paper develops Neural Bellman Operators that construct feasible dynamic policies and qualify their loss under the original economic law. Own-future neural realization and original-owner compilation are joined to reference-centered recursive-risk bounds. We derive an action-aligned supporting certificate whose enclosure width, directional mismatch and constrained stationarity residual control the proposal's original-action gap. Prospective selection then controls rejection or bounded complete work. For convex continuations, nonuniform-law AVaR enclosures attain the deterministic information exponent two; separately quantified randomized lower bounds distinguish sound intervals from point estimation. An exactly solved purchase mechanism turns acquired oblique geometry into a strict policy-loss separation from the best additive and quadratic classes, while identical conventional geometry remains equivalent. The retained continuous-investment evidence favors quadratic least squares and shows that smaller global residual spans can increase rejection. The construction, information and empirical accounts therefore retain distinct economic meanings.
\end{abstract}
\begin{keyword}
\kwd{Dynamic programming}\kwd{Neural approximation}\kwd{Recursive risk}
\kwd{Policy qualification}\kwd{Computational economics}
\end{keyword}
\end{frontmatter}
\section{Introduction}\label{sec:introduction65}
A learned continuation is economically useful when it supports feasible dynamic decisions whose loss can be assessed under the original economic law. Neural Bellman Operators (NBOs) join three objects: an attained continuation, the feasible actions it proposes, and the complete policy that is implemented after numerical qualification. Their common comparison is
\[
 \underline V\le V^*\le J^{\widehat\pi}\le\overline J.
\]
The lower inequality covers every original feasible action. The upper inequality includes observation, action realization and future control. A guarantee for an approximating objective is not a substitute for this comparison.

This paper develops that construction and its decision-relevant resource account. Each dated NBO is formed against its own already attained future, using primitive state and action covers, integration allowances and feasible original owners. The affine--ReLU realization does not assume successful unrestricted nonconvex training. Original-owner compilation changes evaluation work while preserving the continuation and the actual feasible policy. The retained joint-domain, common-readout, horizon and vector extensions keep their original hypotheses and full statements.

For recursive monetary risk, the policy account centers the original conditional dual at the stored continuation. The penalty for changing a maximizing risk density remains part of the comparison. State-dependent lower and upper potentials then compare the complete policy with the original optimum. In a continuous investment economy, original primitive curvature supplies global continuation supports and all-action quadratic lower witnesses. Full-law envelopes bound the implemented controller from above. This makes the qualification constructive without replacing the continuous innovation law by an empirical one.

The acquisition problem is more demanding than certificate validity. Our retained matched experiment finds that every certificate-oriented readout improves its validated residual span, but aggregate first-proposal rejections increase from 644 to 757. Quadratic least squares remains the complete-work winner. We therefore derive an action-aligned certificate that separates the original-objective enclosure width, a supporting-slope mismatch and the surrogate's constrained stationarity error. A supporting slope is not assumed to be an exact derivative of the optimal continuation. Its verified quadratic minorant directly bounds the proposal's all-action gap. Selection over a prospectively fixed menu then has a population rejection guarantee, or a complete-work guarantee when work itself is the selected bounded loss. These results explain which training objective is entitled to predict which service outcome.

We also make the economic role of acquired geometry explicit. A two-date purchase problem uses the original oblique continuation $1+(x_1+x_2-1)_+$. Under the stated uniform state and posted-price laws, expected policy regret equals one sixth of the squared continuation error. The best additive and total-degree-two continuations incur exact average regrets $1/540$ and $1/5400$, whereas an oblique ReLU is exact. The identical conventional hinge and a two-triangle representation are also exact. Thus the example establishes a decision-theoretic value of the geometry, not a distinction between mathematically identical encodings. Two independently implemented rational acquisition calculations reproduce the optimal coefficients and policy losses. This mechanism supplements, rather than reinterprets, the investment evidence.

Information resolution supplies a separate constraint. For convex Lipschitz continuations, we extend the deterministic AVaR information bound and constructive enclosure to known nonuniform densities bounded above and away from zero. At each fixed law and risk level the deterministic error exponent remains two. Pathwise-sound randomized intervals retain that obstruction, whereas ordinary randomized point estimation requires a different argument. A fixed distribution of convex tail perturbations gives an explicit root-mean-square lower bound of order $N^{-5/2}$; we do not identify that exponent as an optimal randomized rate. Dated precision allocation connects the constructive enclosures to the charged Bellman query account while retaining state, action and tensor-product costs.

\paragraph{Related work.}
Lipschitz extension, distance transforms, approximate policy iteration and safe improvement are established tools \citep{mcshane1934,felzenszwalb2012,munos2003,pirotta2013,scherrer2015,yang2020convex}. Risk-averse dynamic programming and feature-based risk learning likewise have an established theory \citep{ruszczynski2010risk,ruszczynski2014erratum,ruszczynskiyang2025}. Derivative information in numerical dynamic programming is developed by \citet{CaiJudd2015}; decision-aware prediction and surrogate calibration are studied by \citet{ElmachtoubGrigas2022,LiuGrigas2021}. HJB verification also connects residual or derivative-sensitive approximation bounds to controller performance \citep{JonesPeet2026,Liu2026HJB}. We do not claim the first link between approximation and control loss. The present support certificate retains the original discrete-date risk law, imperfect feasible witnesses, the full original action set and the implemented policy's numerical obligations. The information arguments use classical deterministic and randomized lower-bound methods \citep{Novak2015Integration97}, specialized here to the stated convex AVaR class.

\paragraph{Evidence and organization.}
The article presents the own-future construction, original-law risk qualification, continuous investment implementation, action-aligned acquisition, information resolution, decision-theoretic geometry and the two distinct investment catalogues. The frozen population comparison and the subsequent acquisition diagnostics are not pooled, and their statistical intervals are not transferred between refitted actors. Acquisition and whole-workload costs enter the purchase account with their correct units; the investment price exercise remains conditional on its declared prices and timing envelope. The Technical Supplement proves the active results. The Development Appendix retains the complete prior expectation, signed-risk, information and acquisition results together with the original broader applications and numerical records. A content map records every relocation; no mathematical result or adverse observation is removed.

\section{The economic operator and its policy account}\label{sec:core49}
Consider dates $t=0,\ldots,T$ and an invariant compact state domain $X\subset\R^d$. At date $t$, the action belongs to a nonempty compact set $A_t(x)$. The transition is $F_t(x,a,Z_{t+1})$, the current cost is $c_t(x,a)$, and the terminal cost is $g(x)$. Write $B_0=1$ and $B_t=\prod_{j<t}\beta_j$, with $\beta_j>0$. An implemented policy includes its observation rule, action selection and feasible numerical realization, not just its stored parameters.

For a continuation $v$, define
\begin{equation}\label{eq:operator49}
 Q_t(v;x,a)=c_t(x,a)+\beta_t\rho_t(v(F_t(x,a,Z))),\qquad
 \mathcal T_tv(x)=\inf_{a\in A_t(x)}Q_t(v;x,a).
\end{equation}
Throughout the deterministic comparison, the evaluations $\rho_t$ are monotone and cash invariant on a common domain containing all functions and constant translates used below. They are therefore nonexpansive in the supremum norm. Continuous objectives, effective compact feasible correspondences and measurable minimizers, or measurable arbitrarily near-minimizers with their error charged, are assumed. The stipulated conditional recursion gives Bellman comparison against all admissible adapted controls. In the earlier expected-cost experiments $\rho_t$ is conditional expectation, costs are bounded and innovations have their original continuous uniform law. Section~\ref{sec:continuous95} evaluates the same investment technology under an explicitly nested AVaR preference and retains the original continuous innovation law. The direct-cost inference in Section~\ref{sec:decisions49} is specific to expectation; it is not automatically an inference theorem for nonlinear recursive utility or games.

A Neural Bellman Operator is a construction and evaluation procedure for \eqref{eq:operator49}. It forms targets using its \emph{own} fitted $f_{t+1}$, returns a neural continuation and a feasible policy, and records an economic error allowance. It does not receive the optimal continuation as a training label. Arbitrary trained networks can be evaluated by the first result below. The subsequent constructive backend specifies how a network and its actor are obtained without assuming successful nonconvex optimization.

\subsection{Action-dependent errors and complete-policy loss}
Let $H(a)$ be a true current cost and $\widehat H(a)$ its continuation approximation. If $\widehat a$ is $\delta$-optimal for $\widehat H$ on the same feasible action set, then
\begin{equation}\label{eq:center49}
 H(\widehat a)-\inf_aH(a)\le\delta+
 \operatorname{osc}_{a}\{\widehat H(a)-H(a)\}.
\end{equation}
Indeed, compare $\widehat a$ with an arbitrarily near-optimal true action and cancel the approximate objectives. The infimum limit gives the result. Additive continuation levels cancel; action-dependent discrepancies do not. Under a change in future policies, the continuation must be transported to the new economic problem before this inequality licenses reuse. The full development edition retains the original quantitative transport, recursive-domain and economic-application results. No transport bound is presumed valid after an unaccounted change of constraints or preferences.

\begin{proposition}[One-sided policy sandwich]\label{prop:sandwich49}
Suppose, uniformly in state,
\[
 f_t-\mathcal T_tf_{t+1}\le u_t,\qquad
 \mathcal T_t^\pi f_{t+1}-f_t\le v_t,
 \qquad \ell_T\le f_T-g\le u_T.
\]
Then the actual feasible policy satisfies
\begin{equation}\label{eq:sandwich49}
 0\le J_0^\pi-V_0^*\le
 \sum_{t<T}B_t(u_t+v_t)+B_T(u_T-\ell_T).
\end{equation}
The bound is invariant to date-specific constant shifts of the fitted continuations, provided the associated residual accounts are shifted consistently.
\end{proposition}
\noindent The proof is in Supplement, Section~\ref{supp:proof-core49-1}.


A sharper valid residual account for the same continuation and the same deployed actor sharpens knowledge of the policy. It does not retrain that policy or change its realized cost. Conversely, subtracting two upper bounds in \eqref{eq:sandwich49} does not identify the difference in actual costs. These distinctions govern both the historical comparisons and the new experiments.

\section{Constructing and compiling a feasible neural policy}\label{sec:construction49}
Assume the primitive objective has, for feasible state--action pairs, state modulus $A_t^Q=C_{x,t}+\beta_tF_{x,t}L_{t+1}$ and action modulus $D_t=C_{a,t}+\beta_tF_{a,t}L_{t+1}$. Let
\[
 d_H(A_t(x),A_t(z))\le K_t^A\|x-z\|_1,
 \qquad L_t=A_t^Q+D_tK_t^A,\qquad L_T=L_g.
\]
The Hausdorff condition and the primitive moduli imply that $\mathcal T_tf_{t+1}$ is $L_t$-Lipschitz whenever $f_{t+1}$ is $L_{t+1}$-Lipschitz. All moduli and numerical integration procedures are supplied by the economic primitives, not fitted from successful outcomes.

At state nodes $x_{t,i}$ with covering radius $h_t$, form effective feasible $k_t$-nets. Evaluate their own-future objectives with certified error at most $e_t$. Minimize the computed labels and retain the corresponding \emph{original} action $a_{t,i}$. A Borel repair $\Gamma_{t,i}(x)$ is feasible at $x$ and satisfies
\begin{equation}\label{eq:repair49}
 \|\Gamma_{t,i}(x)-a_{t,i}\|_1\le K_t^A\|x-x_{t,i}\|_1+\zeta_t.
\end{equation}
The allowance $\zeta_t$ concerns displacement; feasibility itself may not fail. Define
\begin{equation}\label{eq:cone49}
 f_t(x)=\min_i\{y_{t,i}+L_t\|x-x_{t,i}\|_1\},\quad
 i_t(x)=\min\arg\min_i\{y_{t,i}+L_t\|x-x_{t,i}\|_1\}.
\end{equation}
Return $\pi_t(x)=\Gamma_{t,i_t(x)}(x)$. Terminal labels approximate $g$ with error $e_T$ and use the same cone construction with modulus $L_g$.

\begin{theorem}[Own-future construction with feasible witnesses]\label{thm:construct49}
The functions in \eqref{eq:cone49} have exact affine--ReLU realizations and Lipschitz moduli $L_t$ independent of label error. The returned policy is feasible and measurable. It satisfies
\begin{align}\label{eq:enclosures49}
 -e_t\le f_t-\mathcal T_tf_{t+1}&\le D_tk_t+e_t+2L_th_t,\\
 \mathcal T_t^\pi f_{t+1}-f_t&\le e_t+D_t\zeta_t,\notag
\end{align}
and hence
\begin{equation}\label{eq:constructionloss49}
 J_0^\pi-V_0^*\le
 \sum_{t<T}B_t(D_tk_t+2e_t+2L_th_t+D_t\zeta_t)
 +B_T(2e_T+2L_gh_T).
\end{equation}
With effective covers, repairs and certified queries at the prescribed accuracy, allocating each nonnegative term a share of a positive economic tolerance gives a terminating finite construction. A fixed numerical precision or a capped finite experimental ladder need not contain that allocation.
\end{theorem}
\noindent The proof is in Supplement, Section~\ref{supp:proof-core49-2}.


The Lipschitz-extension principle is classical \citep{mcshane1934}. A native min-plus envelope and its exact ReLU circuit, using the same labels, witnesses and ties, are \emph{one method with two encodings}. They return the same function and policy. Neither Theorem~\ref{thm:construct49} nor a faster representation proves a uniquely neural capability or a training-success probability for Adam or L-BFGS. The mathematical increment is the complete feasible, own-future construction and implementation account; the empirical question is its work and policy value relative to conventional generators.

\subsection{Original-owner compilation on a tensor cover}
On a tensor cover $\mathcal X=\prod_{j=1}^d\{x_{j,0},\ldots,x_{j,n_j}\}$, put $S=\prod_j(n_j+1)$. For arbitrary labels $y_i$, including inconsistent inexact labels, compute at each grid node $v$
\[
 (z_v,o_v)=\min_{\rm lex,i}(y_i+L\|v-x_i\|_1,i).
\]
The index $o_v$ is the original label owner, not the transformed node $v$.

\begin{theorem}[Continuous original-owner identity]\label{thm:compile49}
For a closed tensor cell $C$ with corners $\mathcal V(C)$, and $x\in C$,
\begin{align}\label{eq:compile49}
 f(x)&=\min_{v\in\mathcal V(C)}(z_v+L\|x-v\|_1),\\
 (f(x),i_*(x))&=\min_{\rm lex,v\in\mathcal V(C)}
 (y_{o_v}+L\|x-x_{o_v}\|_1,o_v).\notag
\end{align}
Coordinatewise forward and backward relaxations compute $(z,o)$ in $O(dS)$ additions and comparisons and $O(S)$ storage. After locating the cell, a query uses $O(d2^d)$ arithmetic operations. Exact rational comparisons preserve all ties, including $L=0$. Using the original owner's repaired action preserves the uncompiled policy, its cost and its certificate exactly.
\end{theorem}
\noindent The proof is in Supplement, Section~\ref{supp:proof-core49-3}.


The nodal algorithm is the classical distance transform \citep{felzenszwalb2012}; preserving the original action through continuous off-grid queries is the policy obligation. If there are $J_{t,i}$ node actions and $M_t$ integration representatives, write $Q_t=M_t\sum_iJ_{t,i}$. Literal dense cone evaluation costs $O(\sum_t dS_{t+1}Q_t)$ in representation work. Compilation replaces this by $O(\sum_t[dS_{t+1}+Q_td2^d])$, plus cell location, primitive evaluation, integration, feasible repair, certificate and output. Rational bit lengths and state boxes crossing many cells incur additional costs. There is no dimension-free complexity assertion.

\subsection{Acquired states and conventional actors}
A $b$-bit coordinate sensor reports a dyadic cell and its midpoint $\widehat x$. On the unit cube, $\|x-\widehat x\|_1\le r(b)=d2^{-b-1}$. For the investment capacity $\kappa(x)$ below, clip the original action to the smallest capacity throughout that cell and round down to an action quantum $\Delta_t$. This is feasible at every state in the cell. Relative to an exact true-state repair, its displacement is at most $2K_t^Ar(b)+\Delta_t$. Transporting the witness cone from $\widehat x$ to $x$ adds at most $2L_tr(b)$, plus any certified selector objective error $\nu_t$. Thus the implemented-policy account adds
\begin{equation}\label{eq:acquire49}
 \sum_{t<T}B_t\{(2L_t+2D_tK_t^A)r(b)+D_t\Delta_t+\nu_t\}.
\end{equation}
No continuity of the discrete owner selector is assumed. At uncertain numerical comparisons the simulator retains all possible actors, rather than choosing a favorable one.

For conventional FVI the continuation is multilinear on its realized tensor grid. Its actual modulus is bounded by the largest coordinate edge slope. A nearest-node actor uses a feasible node minimizer and true-state repair. Let $h_t$ be the largest half-cell $\ell^1$ radius, $L_t^{\rm gr}$ the Bellman graph modulus, and $e_t$ the node query allowance. Its optimal residual upper bound is $D_tk_t+e_t+L_t^{\rm gr}h_t$. Its selected-policy upper bound is $e_t+E_t$, where
\[
 E_t=\sup_x\{y_{i(x)}+L_t^{\rm gr}\|x-x_{i(x)}\|_1-f_t(x)\}.
\]
Partitioning cells by nearest-node boundaries makes the expression separately affine in each coordinate. Its maximum is therefore at a partition vertex; the implementation encloses all such vertices. Both uniform and curvature-refined FVI receive this sharper one-sided account, not an artificially doubled baseline. Their acquired-state allowance is derived separately in \eqref{eq:fvisensor49}.

\paragraph{Primitive accuracy and resource growth.}
The complete primitive investment construction, its original-optimum accuracy theorem, and the dated horizon allocation are retained in Appendix Sections~\ref{sec:accuracy62} and~\ref{sec:complexity86}. Their state, action, innovation, and bit-growth costs remain explicit. The continuous recursive-risk investment analysis below supplies the corresponding original-law qualification used in the current comparison.


\section{Reference-centered risk kernels and complete-policy qualification}\label{sec:centeredrisk92}
A datewise worst-case error account can charge incompatible states at successive dates. Under nonlinear recursive preferences, replacing that account by an expectation under the physical transition law is not valid. We instead retain the conditional risk evaluation and center its dual representation at the actual stored NBO continuation. The resulting state-dependent certificate has an original-optimum interpretation and a finite occupation representation. It extends, rather than replaces, the signed qualification of Section~\ref{sec:risk91}.

\subsection{The centered operator}
Retain the original augmented economic state, feasible actions, nonnegative discounts, and bounded complete-policy recursions. At each $(t,x,a)$ suppose the original conditional evaluation admits the representation
\begin{equation}\label{eq:dualrisk92}
 \rho_{t,x,a}(v)=\sup_{q\in\mathcal Q_{t,x,a}}
 \left\{\int v(y)q(dy)-\psi_{t,x,a}(q)\right\}.
\end{equation}
Every $q$ is a probability measure on the original successor state space. The representation is assumed on a common bounded-function domain containing all continuations below, with the measurability and selection conditions needed for the original Bellman recursions. This is a hypothesis about the original preference primitive, not permission to replace a static risk objective by a nested one. Conditional risk duality and risk-averse dynamic programming are established principles \citep{ruszczynski2010risk,ruszczynski2014erratum}. For a nonconvex or non-cash-invariant aggregator lacking \eqref{eq:dualrisk92}, the earlier order and modulus results remain applicable.

Fix one qualified parameter $\theta$ and suppress its subscript. Its entire sequence of stored continuations is $f=(f_t)_{t=0}^T$, and its supplied complete feasible policy is $\pi$. Define, using the same $f_{t+1}$ in every term,
\begin{align}\label{eq:centered-def92}
 r_t(x,a)&=c_t(x,a)+\beta_t\rho_{t,x,a}(f_{t+1})-f_t(x),\notag\\
 k_t^f(x,a,q)&=\psi_{t,x,a}(q)+\rho_{t,x,a}(f_{t+1})
                         -\int f_{t+1}\,dq,\notag\\
 \mathcal E_{t,x,a}^f(z)&=\sup_{q\in\mathcal Q_{t,x,a}}
                   \left\{\int z\,dq-k_t^f(x,a,q)\right\}.
\end{align}
The penalty is nonnegative and has infimum zero, whether or not the supremum in \eqref{eq:dualrisk92} is attained. Direct substitution gives
\begin{equation}\label{eq:exact-center92}
 \mathcal E_{t,x,a}^f(z)
 =\rho_{t,x,a}(f_{t+1}+z)-\rho_{t,x,a}(f_{t+1}).
\end{equation}
In particular, $\mathcal E^f(0)=0$, and $\mathcal E^f$ is monotone, cash invariant, and nonexpansive in the uniform norm. Discarding $k^f$ is an upper relaxation, not this identity. Freezing one risk-maximizing kernel at $f$ is in general not even an upper relaxation.

\begin{theorem}[State-dependent original-optimum certificate]\label{thm:centeredrisk92}
Suppose verified functions satisfy $r_t^-(x,a)\le r_t(x,a)$ for every original feasible action and $r_t(x,\pi_t(x))\le r_t^{\pi,+}(x)$. Let bounded measurable $L_t,U_t$ satisfy
\begin{align}\label{eq:potential-cert92}
 L_T&\le g-f_T\le U_T,\notag\\
 L_t(x)&\le\inf_{a\in A_t(x)}
       \{r_t^-(x,a)+\beta_t\mathcal E_{t,x,a}^f(L_{t+1})\},\notag\\
 U_t(x)&\ge r_t^{\pi,+}(x)
                +\beta_t\mathcal E_{t,x,\pi_t(x)}^f(U_{t+1}).
\end{align}
Then, at every initial state,
\begin{equation}\label{eq:potential-loss92}
 f_0+L_0\le V_0^*\le V_0^\pi\le f_0+U_0,
 \qquad 0\le V_0^\pi-V_0^*\le U_0-L_0.
\end{equation}
The same statements hold at every intervening date and state. They are uniform over a qualified parameter class whenever all premises and the reported bound hold for each member with its own common continuation sequence and complete policy. The certificate permits endogenous states and actions changing at every date; it does not require a finite state space or independent deployment contexts.

With exact residuals and terminal discrepancies, equality recursions in \eqref{eq:potential-cert92} give $L_t=V_t^*-f_t$ and $U_t=V_t^\pi-f_t$. With scalar bounds from \eqref{eq:signed-input91}, their corresponding constant lower and upper sequences are feasible certificates here. Thus the exact state-dependent recursion is no looser than the signed scalar account. A strictly smaller bound can result without changing a fit, a feasible action, or an economic observation.
\end{theorem}

The last assertion concerns the best exact recursion or a verified improvement upon the constant certificates. It is not a claim that every approximately computed $L,U$ improves the scalar account. In implementation, the minimum of valid policy supersolutions and the maximum of valid optimum subsolutions retain validity, so an inferior new bound need not displace the existing one.

The lower inequality requires every original action. A discrete action search without a verified original-menu lower bound cannot establish it. For a fixed action, a feasible risk kernel supplies a lower bound on $\mathcal E^f(L)$; an upper bound requires control of the entire risk envelope. Numerical integration, feasible repairs, residual errors and penalty errors belong in these inequalities. If penalty enclosures are used, lower penalty endpoints enter an upper certificate and upper endpoints enter a lower witness. Approximate dual optimizers are accepted only through these independently checked inequalities.


\section{Continuous-law qualification in the investment model}\label{sec:continuous95}
The risk theorem becomes a numerical method only when its inequalities can be constructed from economic primitives. We now supply this construction for the original investment technology with a nested tail-risk preference. States and innovations remain continuous, every feasible action is covered, and the implemented terminal controller rather than a hypothetical exact optimizer determines policy cost. This application joins continuation construction, the original-optimum comparison and a representation comparison on one problem. It does not replace the general NBO construction or the other economic applications.

\subsection{Primitives and the object being approximated}
Let $d\geq2$ be even, $X=[0,1]^d$, and interpret subscripts cyclically. At dates $t=0,1$, the feasible investment is $0\leq a\leq A(x)=1/8+\bar x/8$, where $\bar x=d^{-1}\sum_jx_j$. Write $b_j=1/2$ for odd $j$ and $b_j=1/4$ for even $j$, and $s_j=(-1)^{j+1}$. The original transition is
\begin{equation}\label{eq:investment95}
 F_j(x,a,z)=\frac1{16}+\frac{x_j}{2}+\frac{x_{j+1}}8
       +\frac{x_j(1-x_{j+1})}{16}+b_ja+s_jz,
 \qquad z\sim\operatorname{Unif}[-1/32,1/32].
\end{equation}
Define
\begin{align}\label{eq:cost95}
 S(x)&=\frac2d\sum_j(x_j-5/8)^2+
       \frac1{4d}\sum_j(x_j-x_{j+1})^2+2(1/2-2\bar x)_+^2,\\
 g(x)&=S(x)+\frac2d\sum_j(x_j-5/8)^2,\qquad h(a)=a^2+4a^4.\nonumber
\end{align}
We use $\beta=15/16$ and, separately, $\alpha\in\{1/2,3/4\}$. Average value at risk is the upper-tail cost functional
\[
 \rho_\alpha(v)=\sup\left\{\int vq\,dP:
             0\leq q\leq(1-\alpha)^{-1},\ \int q\,dP=1\right\}.
\]
Here $P$ is the original uniform innovation law. The objective is recursively defined by $V_2=g$ and
\[
 V_t(x)=S(x)+\inf_{0\leq a\leq A(x)}
       \{h(a)+\beta\rho_\alpha(V_{t+1}(F(x,a,\cdot)))\}.
\]
It is not the static tail risk of total realized cost. A complete policy value replaces each infimum by the action prescribed by that policy, with the same recursive risk primitive.

\paragraph{Assumptions for the continuous construction.}
For this subsection, the primitives are exactly \eqref{eq:investment95}--\eqref{eq:cost95}; $\alpha$ is fixed before construction; the fitted continuation $f$ is a finite convex function on $X$; and every numerical lower or upper value is an outward enclosure. Feasible tail densities and feasible actions are retained explicitly. Terminal selection is the measurable, dyadically acquired controller described below, with its fallback included in the error bound. These concrete conditions imply compact feasible graphs, bounded costs, measurable minimizers, and the conditional-risk dual used in Theorem~\ref{thm:centeredrisk92}. They do not assert that every recursive preference or every trained network has this convex structure.

\begin{proposition}[Primitive curvature and state supports]\label{prop:curvature95}
For either risk level and every even $d$, the one-date objective
$Q_1(x,a)=S(x)+h(a)+\beta\rho_\alpha(g(F(x,a,\cdot)))$
is jointly convex on $X\times[0,1/4]$, with curvature bounded below by
\begin{equation}\label{eq:curvature95}
 \frac{4-5\beta/4}{d}\|v\|_2^2+2w^2
\end{equation}
in a direction $(v,w)$. Its value $V_1$ is convex on $X$ and satisfies
\begin{equation}\label{eq:lipschitz95}
 |V_1(x)-V_1(y)|\leq\frac{15}{d}\|x-y\|_1.
\end{equation}
Let $W$ be the objective evaluated under any fixed feasible risk density, at an anchor $(x_0,a_0)$, and let $(p_x,p_a)$ be its gradient there. For any $\lambda,\nu\geq0$, set $e=p_a+\lambda-\nu$. Then the following is a global affine minorant of $V_1$:
\begin{align}\label{eq:affine95}
 \ell(y)={}&W+\lambda(a_0-A(x_0))-\nu a_0-\frac{e^2}{4}\nonumber\\
 &+\left(p_x-\frac{\lambda}{8d}\boldsymbol{1}\right)^{\!\top}(y-x_0).
\end{align}
Neither optimal multipliers nor an exact optimizing action are required.
\end{proposition}

The curvature is obtained from the original nonlinear technology, not imposed on an interpolated value grid. The negative second-order term from $x_j(1-x_{j+1})$ is bounded using $|\partial_jg|\leq10/d$; the stage curvature dominates it. The allowance $e^2/4$ pays for imperfect stationarity. The proof in the Technical Supplement also establishes the global state modulus in \eqref{eq:lipschitz95}, including the state-dependent action constraint.

\subsection{An all-action certificate under the original innovation law}
Partition the innovation interval into $n$ equal bins $B_i$, with midpoints $z_i$ and probabilities $p_i=1/n$. An innovation bin is a set over which an inequality holds, not a replacement point mass in the economic model. Choose nonnegative $q_i\leq p_i/(1-\alpha)$ summing to one. These masses define a feasible density constant on each bin. For a fixed initial state and feasible proposal $a_0$, let $W_0$ be a lower bound on the objective under this density, obtained from globally valid affine supports of the continuation at the successor midpoints. Let $p$ enclose the derivative of the same supporting objective at $a_0$. Then
\begin{equation}\label{eq:actionlower95}
 \underline V_0(x)=\inf_{0\leq u\leq A(x)}
         \{W_0+p(u-a_0)+(u-a_0)^2\}
\end{equation}
is a lower bound on the original Bellman optimum when the supported continuation is $V_1$. The scalar quadratic is minimized over the full original interval. Interval uncertainty in $p$ is covered by taking the smaller minimum at its two endpoints, since the minimum is concave in $p$.

For a convex fitted continuation, midpoint supports give the lower risk witness and the larger endpoint value bounds the continuation throughout each bin. For $V_1$, Proposition~\ref{prop:curvature95} supplies midpoint supports and the allowance $15/(32n)$ covers each half-bin. Upper tail risk is the average of the largest $n(1-\alpha)$ upper bin bounds. This average is the exact risk of a dominating step function, hence an upper bound under the continuous law.

\begin{theorem}[Constructed original-optimum qualification]\label{thm:continuous95}
Suppose the supports and bin bounds above are outwardly verified. Suppose further that the complete terminal controller satisfies
$V_1\leq V_1^{\widehat\pi}\leq V_1+\varepsilon_1$ everywhere on $X$.
For every queried initial state, the lower value \eqref{eq:actionlower95} and the binwise upper bound $\overline J_0^{\widehat\pi}$ satisfy
\begin{equation}\label{eq:originalgap95}
 0\leq J_0^{\widehat\pi}(x)-V_0(x)
       \leq\overline J_0^{\widehat\pi}(x)-\underline V_0(x).
\end{equation}
For any stored convex continuation $f$, the same lower and upper constructions are admissible bounds for the centered-risk comparison of Theorem~\ref{thm:centeredrisk92}. Their validity does not depend on the representation of $f$, on a successful fit, or on the proposal being an exact minimizer.
\end{theorem}

To see the connection, write a verified continuation bracket as $\ell\leq V_1\leq u$ and use errors $L=\ell-f$, $U=u-f$. The centered identity gives
\[
 \rho_\alpha(f)+\mathcal E^f(L)=\rho_\alpha(\ell),\qquad
 \rho_\alpha(f)+\mathcal E^f(U)=\rho_\alpha(u).
\]
Joint evaluation avoids subtracting and then adding two separately rounded copies of $\rho_\alpha(f)$. The bound \eqref{eq:actionlower95} uses one feasible risk density to minorize the required all-action infimum; the policy upper bound controls the full risk envelope. Thus the application constructs sufficient inequalities for the original theorem. It does not claim that a non-neural continuation would receive a different validity theorem, or that computing the original-value witness is free.

The terminal controller observes each state coordinate rounded down to $40$ binary places. It accepts its floating proposal only after an outward all-action gap of at most $2^{-21}$ at that observed state; otherwise it uses rational bisection. The original-state transfer is less than $30\,2^{-40}$. Consequently $\varepsilon_1=2^{-20}$ covers the complete controller, including observation and fallback. Policy evaluation encloses this controller at every possible successor, not only at the midpoint states used by the integrator.

\subsection{Residual covers and strict continuum screening}
A cellwise residual bracket need not solve another action problem at every state. On each rectangular state cell, convex interpolation of the vertex upper values dominates $V_1$; each affine function \eqref{eq:affine95} minorizes it globally. The same interpolation and subgradient construction applies to convex $f$. Comparing an affine minorant of one function with the convex interpolation upper bound of the other at every vertex proves a bound for $V_1-f$ throughout the cell. Taking the best available anchor and then the worst cell gives a whole-domain residual interval. The $32\times32$ cover below therefore qualifies the continuous square, rather than converting its $1,089$ nodes into empirical generalization observations.

Let $[l_i,u_i]$ enclose an original successor value over every point of bin $B_i$. If $\tau$ is the $k$th largest $l_i$, $k=n(1-\alpha)$, then
\begin{equation}\label{eq:strict95}
 u_i<\tau
\end{equation}
excludes all risk mass from bin $i$ in an upper-tail maximizer. At least $1-\alpha$ of the original probability mass has value at least $\tau$; strict separation makes any positive mass on the excluded bin suboptimal. Equality is retained. These bounds must be available before screening. The screening calculation below pays for their construction and for the test itself, and compares identical upper-tail values with and without screening. This continuum-bin result is separate from enumerating all extreme kernels of a finite risk envelope.


\section{Action-aligned acquisition and complete-policy work}\label{sec:action99}
A continuation is acquired to make a decision, not to minimize an approximation statistic in isolation. The retained investment experiment makes the distinction operational: every certificate-oriented fit reduces its validated residual span, whereas aggregate first-proposal rejections increase. We now construct an acquisition target in the units of the admission inequality. The construction uses the same original action set and risk law as the NBO policy account.

\subsection{A computable directional certificate}
Fix a state and date, and suppress their indices. Let $A\subset\mathbb R^k$ be nonempty, compact and convex, let $Q$ be the original action objective, and let $a\in A$ be the actually realized proposal. A verified supporting calculation supplies $L,U\in\mathbb R$, $s\in\mathbb R^k$ and $\mu>0$ such that
\begin{equation}\label{eq:support99}
 Q(a)\le U,\qquad
 Q(b)\ge L+s^{\mathsf T}(b-a)+\frac\mu2\|b-a\|_2^2
 \quad\text{for every }b\in A.
\end{equation}
These inequalities need not identify $s$ as a derivative of $Q$. In particular, an imperfect risk witness or continuation minorant can lower $L$ while retaining a useful supporting slope. Put $w=U-L$ and
\begin{equation}\label{eq:gamma99}
 \Gamma(a)=w-\min_{b\in A}
       \left\{s^{\mathsf T}(b-a)+\frac\mu2\|b-a\|_2^2\right\}.
\end{equation}
A certified lower bound on the displayed minimum gives a valid, possibly larger, $\Gamma$. An unverified optimizer status does not.

\begin{theorem}[Directional certificate for an NBO proposal]\label{thm:action99}
Under \eqref{eq:support99},
\[
 0\le Q(a)-\inf_{b\in A}Q(b)\le\Gamma(a).
\]
Let $\widehat s$ be a supporting slope of the surrogate used to propose $a$, let $N_A(a)=\{n:n^{\mathsf T}(b-a)\le0\ (b\in A)\}$, and suppose
$\|s-\widehat s\|_2\le d$. If a supplied $n\in N_A(a)$ satisfies
$\|\widehat s+n\|_2\le r$, then
\begin{equation}\label{eq:direction99}
 \Gamma(a)\le w+\frac{(d+r)^2}{2\mu}.
\end{equation}
For $A=[0,A_0]$, the exact minimum in \eqref{eq:gamma99} is attained at
$b=\operatorname{proj}_{[0,A_0]}(a-s/\mu)$.
\end{theorem}

The proof, including boundary actions, is in Supplement~\ref{supp:action99}. Equation~\eqref{eq:direction99} separates three trainable or computable quantities: the original-objective enclosure width $w$, the directional mismatch $d$, and the surrogate stationarity residual $r$. The residual span of the continuation is not one of these quantities. A readout can improve that span and worsen the supporting slope at a nearly binding action. Conversely, constant changes in a continuation need not affect any of them after consistent centering.

The supporting calculation is available under the original recursive-risk primitives. Suppose $F(b,\omega)=F(a,\omega)+G(\omega)(b-a)$, the stage cost is $\mu$-strongly convex, and the risk dual has an action-independent feasible witness $q$ on the original innovation space. If
\[
 V(F(b,\omega))\ge \ell(\omega)+v(\omega)^{\mathsf T}G(\omega)(b-a),
\]
then \eqref{eq:support99} holds with
\begin{equation}\label{eq:lift99}
 L=c(a)+\beta\left\{\int\ell\,dq-\psi(q)\right\},\qquad
 s=\nabla c(a)+\beta\int G^{\mathsf T}v\,dq,
\end{equation}
and any valid upper bound $U$ for the original risk objective. Integration and penalty errors must be enclosed in these quantities. The witness must remain feasible for every action being minorized; a witness feasible only at $a$ is insufficient when the risk set itself depends on the action.

In the investment economy, $G=(1/2,1/4)^{\mathsf T}$ at $d=2$, the innovation law and AVaR cap are action independent, and $h''\ge2$. Thus $\mu=2$ gives exactly the all-action quadratic of \eqref{eq:actionlower95}. The added result is the explicit link from an acquired directional support to its admission gap, rather than a new economic model or an assumption that an acquired value support is an exact derivative.

\subsection{Which training loss controls rejection?}
Fix a finite menu $\mathcal H$ of $M$ complete candidate services before validation. Each candidate specifies its proposal, rounding, certificate algorithm and fallback. On a context $x$, its first proposal is accepted when the valid complete-policy certificate $\Gamma_h(x)$ is at most $\tau>0$. For a one-date problem, \eqref{eq:gamma99} is such a certificate. At multiple dates it must additionally include the certified future-policy allowance. Define
\begin{equation}\label{eq:loss99}
 R_h(x)=\boldsymbol 1\{\Gamma_h(x)>\tau\},\qquad
 \ell_h(x)=\min\{1,\Gamma_h(x)/\tau\}.
\end{equation}
Then $R_h\le\ell_h$. Equality at the admission boundary is accepted. A valid fallback completes every rejected context; discarding it would change the service and invalidate its resource comparison.

\begin{theorem}[Prospective rejection and work guarantees]\label{thm:selection99}
Condition on all acquisition information that fixes $\mathcal H$. Let $X_1,\ldots,X_n$ be independent draws from a declared deployment-context law, independent of that information. For $0<\delta<1$, set
\[
 e_n=\sqrt{\frac{\log(2M/\delta)}{2n}}.
\]
An empirical minimizer $\widehat h$ of $R_h$ satisfies, with conditional probability at least $1-\delta$,
\begin{equation}\label{eq:rejection99}
 \mathbb E R_{\widehat h}(X)\le\min_{h\in\mathcal H}\mathbb E R_h(X)+2e_n.
\end{equation}
For an empirical minimizer of $\ell_h$, the same event instead gives
$\mathbb E R_{\widehat h}\le\min_h\mathbb E\ell_h+2e_n$.
If a verified resource bound is $W_h(x)\le C_0+C_1R_h(x)$, with common nonnegative constants, the rejection selector satisfies $\mathbb E W_{\widehat h}\le C_0+C_1(\min_h\mathbb E R_h+2e_n)$, and the surrogate selector has the analogous bound with $\ell_h$. More generally, directly minimizing empirical $W_h/W_{\max}\in[0,1]$ yields
\[
 \mathbb E W_{\widehat h}\le\min_h\mathbb E W_h+2W_{\max}e_n.
\]
All selection and rejected-candidate validation costs are additional one-time acquisition charges.
\end{theorem}

The work statement means $\mathbb E W_{\widehat h}\le C_0+C_1(\min_h\mathbb E R_h+2e_n)$ for the rejection selector, and analogously for $\ell$. A measured maximum is not a verified $W_{\max}$, and a rejection-optimal selector need not minimize runtime when fallback costs differ across contexts. On the finite dyadic domain used by the investment service, a finite fallback can supply a worst-case operation cap; its size, rather than a favorable observed clock, enters this theorem. Conditional concentration allows data-dependent acquisition, but does not allow the menu or certificate rule to be changed after looking at the validation sample.

Theorems~\ref{thm:action99}--\ref{thm:selection99} provide the missing distinction between optimizing a residual certificate and optimizing an action-relevant quantity. They do not retroactively certify the empirical performance of the R97 residual objective, or supply observations for the separately frozen R98 actors. Decision-aware prediction and surrogate calibration are established ideas \citep{ElmachtoubGrigas2022,LiuGrigas2021}. Here the loss is attached to an outward original-law, all-action Bellman certificate and a complete feasible service. The finite-menu concentration argument is classical; the support-to-admission construction specifies what that argument is permitted to select.

\section{Decision-scale numerical and statistical precision}\label{sec:precision99}
The continuous qualification uses the original terminal-controller allowance and all-action lower witness, not just a small surrogate residual. Convex midpoint and endpoint risk rules supply a width of at most $KR/\{(1-\alpha)n^2\}$ with $2n+1$ continuation queries. Theorem~\ref{thm:nonuniform99} below extends this account beyond the uniform law. The full uniform-rule statement, numerical endpoint charges and tensor extension remain in Section~\ref{sec:precision96} of the Development Appendix.

For the frozen population comparison, every complete service is qualified to the original-optimum allowance $\tau=1/256$. The design uses 4,096 independent initial states at each risk level and twelve simultaneous contrasts, with numerical half-width and sampling allowance totaling $3/8192$. Theorem~\ref{thm:power96} states the exact simultaneous event and its independence requirements. Those population intervals belong to the frozen actors, not to subsequently refitted candidates. The acquisition comparison instead retains the deterministic contrast support $[-\tau,\tau]$.

For a workload of $B$ decisions, the purchase account combines the per-decision policy contrast, $B$ times that contrast, full-workload computation and one-time acquisition. Theorems~\ref{thm:purchase97} and Proposition~\ref{prop:purchase96} retain the complete interval and strict-boundary statements. Historical timing ranges are descriptive inputs to a conditional price exercise. The exact decision mechanism in Section~\ref{sec:decision99} will provide an identified model-specific policy contrast; it will not turn these historical clock ranges into confidence intervals or external prices into estimated welfare weights.

\section{Information resolution under nonuniform laws and randomization}\label{sec:information99}
The information required to qualify a continuation is distinct from the information supplied when its formula is already known. This section studies point-value access to a convex, $K$-Lipschitz continuation on $[-R,R]$, with $K,R>0$. The continuation formula, derivatives and exact Bellman labels are not supplied for free. The original uniform-law result, including its complete proof and constants, is retained as Theorem~\ref{thm:information97} in the Development Appendix. We extend the law and distinguish two nonequivalent meanings of randomized certification.

\subsection{A known nonuniform innovation law}
Let $P$ have known density $p$ satisfying $0<m\le p(z)\le M_p<\infty$ on $[-R,R]$. Fix $0\le\alpha<1$, put $t=1-\alpha$, let $q$ be the $\alpha$-quantile of $P$ (with $q=-R$ at $\alpha=0$), and set $L=R-q$. Let $\mathcal C_{K,R}$ denote the convex $K$-Lipschitz functions on this interval.

\begin{theorem}[Nonuniform-law information and enclosure]\label{thm:nonuniform99}
Every deterministic adaptive procedure using at most $N$ point values and reporting an interval valid for all consistent $f\in\mathcal C_{K,R}$ has an input with reported width at least
\begin{equation}\label{eq:lower99}
 \frac{mKL^2}{12t(N+1)^2}
 \ \ge\ \frac{mKt}{12M_p^2(N+1)^2}.
\end{equation}
For an upper bound, partition $[-R,R]$ into $n$ equal cells. Suppose the original cell masses and conditional cell means are available exactly. The mass-weighted conditional-mean rule and the mass-weighted endpoint secant rule bracket $\operatorname{AVaR}_\alpha^P(f)$ and have width at most
\begin{equation}\label{eq:upper99}
 \frac{2M_pKR^2}{t n^2}.
\end{equation}
They use at most $2n+1$ distinct continuation values. Thus the deterministic error-versus-information rate is $N^{-2}$ for each fixed admissible law and fixed $\alpha$. Certified, rather than exact, law moments require their additional numerical allowances.
\end{theorem}

The lower bound uses indistinguishable convex tails under the original density, not a probability transformation that might destroy convexity. The upper bound applies conditional Jensen and secant inequalities to $(f-\eta)_+$ before minimizing over the AVaR threshold $\eta$. The moments are law information and their acquisition cost is explicit. With $m=M_p=1/(2R)$ and $L=2Rt$, both constants reduce to their uniform-law counterparts. Their opposite behavior as $\alpha$ tends to one remains; no tail-uniform constant is identified.

\subsection{Sound random intervals and randomized point estimates}
A randomized algorithm can mean a zero-failure enclosure or a statistical approximation. We call an interval procedure \emph{pathwise sound} if, outside one null set of random tapes independent of the input, it contains the risk of every admissible input consistent with the tape's transcript. This is stronger than requiring coverage $1-\delta$ separately for each input.

\begin{proposition}[Randomized information with explicit quantifiers]\label{prop:random99}
The first lower bound in \eqref{eq:lower99} holds on one fixed quadratic-tail input for every sound random tape of a pathwise-sound interval procedure. It therefore also bounds expected width on that input.

For ordinary randomized point-valued procedures under the uniform law, let $A(f)$ use at most $N\ge1$ point values on every realization. Then
\begin{equation}\label{eq:randomlower99}
 \sup_{f\in\mathcal C_{K,R}}
 \left\{\mathbb E\left|A(f)-\operatorname{AVaR}_\alpha^P(f)\right|^2\right\}^{1/2}
 \ge\frac{KR(1-\alpha)}{960\,N^{5/2}}.
\end{equation}
The expectation is over the algorithm's randomization. This is a lower bound, not a claim that $N^{-5/2}$ is the optimal randomized rate on the whole class.
\end{proposition}

For the first assertion, the adverse quadratic is fixed before the random tape; its indistinguishable chord continuation can depend on the realized transcript because soundness is simultaneous over inputs. That reasoning cannot be applied to a merely high-probability interval. The second assertion instead uses a fixed finite distribution of convex perturbations and the conditional variance of unrevealed signs. Its proof supplies all constants and treats adaptive queries. Classical information-based complexity already uses indistinguishable inputs and randomized bump distributions \citep{Novak2015Integration97}; the result here specializes those tools to the convex monotone AVaR tail and states which certification quantifier survives randomization.

\subsection{Dated resolution and the Bellman resource account}
Suppose, on a fixed acquired state--action graph, date $j$ contributes a verified nonnegative allowance $a_j/n_j^2$ to the original-policy tolerance and incurs $c_j n_j$ continuation-query work, with $a_j,c_j>0$. The coefficients include discounting, the appropriate tail-risk factor and the number of original objectives to be evaluated. The continuous allocation problem has value
\begin{equation}\label{eq:allocation99}
 \min_{n_j>0:\,\sum_j a_j/n_j^2\le\varepsilon}
       \sum_j c_jn_j
 =\frac{\left(\sum_j a_j^{1/3}c_j^{2/3}\right)^{3/2}}{\sqrt\varepsilon}.
\end{equation}
Rounding each minimizing $n_j$ upward preserves the tolerance and adds less than $\sum_jc_j$ query work. This is a constructive composition of certified error accounts, not an end-to-end Bellman lower bound. Acquiring the state--action graph, the neural continuation, law moments, exact arithmetic and feasible owners remains additional work. With multiple innovations a direct tensor implementation still has product point counts. The original horizon and vector-control results in Proposition~\ref{prop:horizon95} and the Development Appendix are retained, including their adverse executed scaling evidence.

\section{A decision-theoretic value of acquired oblique geometry}\label{sec:decision99}
An approximation separation matters economically only when it changes an implemented decision. We give such a realization for the original oblique continuation, without altering the retained investment comparison. This is an exactly solved economic mechanism, not an empirical calibration or a timing experiment.

A decision maker observes $x=(x_1,x_2)\in[0,1]^2$ and a posted price $p\in[0,3]$. Purchasing one durable claim yields next-date benefit
\begin{equation}\label{eq:benefit99}
 v(x)=1+(x_1+x_2-1)_+.
\end{equation}
The feasible action is $b\in\{0,1\}$, and total cost is $b\{p-v(x)\}$. This is a two-date Bellman problem with a deterministic terminal payoff and unit discount. A stored continuation $h(x)\in[0,3]$ implements $b_h(x,p)=\boldsymbol1\{p\le h(x)\}$. The original optimum buys at $p\le v(x)$. A price tie has zero loss. Let states be uniform on the square and prices uniform on $[0,3]$, independently.

\begin{theorem}[Exact policy-loss separation]\label{thm:decision99}
For every measurable $h:[0,1]^2\to[0,3]$, the conditional and unconditional policy regrets are
\begin{equation}\label{eq:decisionloss99}
 \mathbb E_p[J^{b_h}(x,p)-J^*(x,p)]
       =\frac{(h(x)-v(x))^2}{6},\qquad
 \mathbb E[J^{b_h}-J^*]=\frac{\|h-v\|_{L^2}^2}{6}.
\end{equation}
Among arbitrary square-integrable additive continuations and among polynomials of total degree at most two, respectively, the optimal expected regrets are
\begin{equation}\label{eq:separation99}
 \frac1{540}\quad\text{and}\quad\frac1{5400}.
\end{equation}
They are attained by the admissible functions
\begin{align}
 h_A(x)&=\frac56+\frac{x_1^2+x_2^2}{2},\label{eq:additive99}\\
 h_Q(x)&=\frac{17}{15}-\frac35(x_1+x_2)
              +\frac12(x_1^2+x_2^2)+\frac65x_1x_2.\label{eq:quadratic99}
\end{align}
A single oblique ReLU with the affine constant in \eqref{eq:benefit99} attains zero regret. A conventional implementation of the identical oblique hinge also attains zero regret and the same decisions.
\end{theorem}

The shift by one is part of the economic construction. It keeps the projection optima inside the posted-price support, so clipping does not secretly enlarge a comparator class. The proof computes the additive conditional-expectation projection and the six-dimensional polynomial Gram system exactly. It establishes a policy comparison against the best member of each specified class, rather than comparing arbitrarily chosen fits. The scalar price integration is the step that turns continuation error into action loss; no such identity is presumed in the investment economy.

The sup-norm comparison also has a decision interpretation. Under the same conditional price law, the smallest worst-state expected regrets for the two classes are $1/96$ and $1/1536$, respectively. These follow from the sharp errors $1/4$ and $1/16$ in Proposition~\ref{prop:ridge97}, with its explicit feasible attaining functions shifted by one. The population optima in \eqref{eq:additive99}--\eqref{eq:quadratic99} are different from the minimax optima. An acquisition objective must specify which context criterion it is intended to serve.

\subsection{Acquisition, conventional geometry and purchase}
The mechanism identifies a value of the oblique continuation geometry, not a value of relabeling its formula as a network. A piecewise-affine conventional representation on the two triangles separated by $x_1+x_2=1$ is also exact. Convex regression with that hinge already in its dictionary contains the same exact continuation. Those controls are included in the exact comparison. Discovering the separating hyperplane from unknown data, fitting a larger dictionary or constructing an adaptive triangulation are different acquisition problems; their costs cannot be inferred from this formula-access calculation.

Suppose the exact oblique service has additional fixed acquisition cost $K\ge0$, incremental per-decision computation cost $c$, and fee $F\ge0$, all expressed in the units of this model's cost. Against the best quadratic continuation, the exact gain from $m$ independent decisions is
\begin{equation}\label{eq:exactpurchase99}
 G_Q(m)=m(1/5400-c)-K-F.
\end{equation}
When $c<1/5400$, the first strictly beneficial workload is
$\lfloor(K+F)/(1/5400-c)\rfloor+1$. This is a nonzero, exactly identified policy benefit with an explicit acquisition threshold. It is not an estimate of the price of computation in the retained investment application. Equal hinge and neural implementations share this gain when their complete costs agree.

\begin{table}[htbp]
\caption{Exact original-law policy loss in the oblique-continuation mechanism}\label{tab:decision99}
\centering\small
\begin{tabular}{@{}p{.52\linewidth}p{.17\linewidth}p{.22\linewidth}@{}}
\toprule
Continuation class or supplied geometry & Optimal mean regret & Information boundary\\\midrule
Arbitrary square-integrable additive functions & $1/540$ & Population projection\\
Total-degree-two polynomials & $1/5400$ & Six exact moments\\
One oblique ReLU and affine constant & $0$ & Supplied hinge\\
Identical conventional oblique hinge & $0$ & Same supplied hinge\\
Piecewise affine on the two separating triangles & $0$ & Same separating line\\
Convex dictionary containing the oblique hinge & $0$ & Contains exact function\\\bottomrule
\end{tabular}
\par\smallskip\parbox{.96\linewidth}{\footnotesize Notes: States are uniform on $[0,1]^2$ and independent prices uniform on $[0,3]$. Loss is evaluated against the original purchase optimum. The last four rows are equivalent exact geometric realizations, not four independent fitted experiments. The quadratic row uses six right-side moments and its known Gram matrix; moment acquisition is not free when the continuation formula is unknown. No timing, empirical calibration or dimension-wide neural dominance is inferred.}
\end{table}

The table is generated from two separately implemented exact-arithmetic calculations. The production path assembles monomial moments and solves normal equations by rational elimination. The independent path integrates polynomials after a triangular change of variables and uses an orthogonal shifted-Legendre basis. It imports neither the production integrator nor its solver. Both reproduce the fitted function, the original-law loss and the range checks. This independently checks acquisition for the stated finite-dimensional mechanism; it does not independently refit the R97 nonlinear dictionary search or complete the separately frozen R98 campaign.


\section{A remotely frozen, decision-scale comparison}\label{sec:study96}
The new comparison follows the complete-service definition of Theorem~\ref{thm:power96}. It is not a relabeling of the R95 actors. Oblique ReLU, adaptive coordinate-spline and projected-quadratic continuations have three nonlinear features and three affine coefficients. Each fitted design has rank six, uses 1,024 common original-risk training labels, and occupies an identical 128-byte executable encoding. Forward selection considers the same number of candidates and uses identical optimization stopping parameters. Every candidate fit, rank rejection and selected feature is retained. Effective design rank is reported separately from nonzero readout count and conditioning. Equal nominal or solver-call budgets do not imply equal fitting time.

The spline knots use training-state quantiles, and all feature selection uses training data only. Calibration uses 128 separate initial states at each risk level and fixes the selected service before confirmation. The source, six binary models, analysis rule and calibration selection are committed remotely before the new evaluation stream can open. Every arm remains in the comparison. This rank-six experiment repairs one specific confound in the preceding comparison; it is not an exhaustive frontier over ranks, labels, adaptive convex regression and computational budgets.

There are 4,096 new independent initial states per risk level, drawn from the explicitly specified uniform 20-bit dyadic distribution. Innovations between decisions remain continuous. All first proposals, rejected proposals, final actions, refinement levels and fallback flags are recorded. An accepted final action satisfies a $1/256$ original-optimum target and a $1/8192$ cost-width target. The comparison margin is $1/1280$, and the predeclared simultaneous half-width is strictly smaller than half that margin. None of these quantities may change after evaluation.

\IfFileExists{tables/outcomes96.tex}{% Source commit: 53a54e1ae0986057e1b7a7c70a12d439a6b4bd14; freeze SHA-256: 99b2cf1b030b24a0ab7a7c8896da52bce5bbde085549e5ec0c5776215cdc83c6
\begin{table}[htbp]
\caption{New-stream complete services at the original-optimum target}\label{tab:services96}
\centering\small
\begin{tabular}{l l r r r r}\toprule
Risk & Service & \shortstack{Mean cost\\midpoint} & \shortstack{Mean numerical\\width} & \shortstack{Maximum\\loss} & \shortstack{Rejected\\proposals}\\
\midrule
0.50 & Online risk-DP & 0.6197369 & 0.0000011 & 0.0000017 & 0 \\
0.50 & Quadratic & 0.6198883 & 0.0000011 & 0.0025887 & 0 \\
0.50 & ReLU & 0.6199340 & 0.0000011 & 0.0039036 & 1136 \\
0.50 & Adaptive spline & 0.6199407 & 0.0000011 & 0.0039033 & 1275 \\
0.75 & Online risk-DP & 0.6308045 & 0.0000012 & 0.0000033 & 0 \\
0.75 & Quadratic & 0.6309583 & 0.0000012 & 0.0029466 & 0 \\
0.75 & ReLU & 0.6309673 & 0.0000012 & 0.0039032 & 1266 \\
0.75 & Adaptive spline & 0.6310478 & 0.0000012 & 0.0039006 & 1073 \\
\bottomrule\end{tabular}
\par\smallskip\parbox{.97\linewidth}{\footnotesize Each row has 4,096 independent initial states, paired across services. Mean cost uncertainty here is numerical, not a sampling interval. Rejected first proposals receive the separately paid conventional fallback; their original actions and loss bounds remain in the raw data. The target distribution is the declared 20-bit dyadic initial law, with continuous innovations.}
\end{table}
\begingroup\small\setlength{\tabcolsep}{3pt}
\begin{longtable}{@{}l p{.24\linewidth} r r r@{}}
\caption{All prespecified paired mean-cost comparisons}\label{tab:pairs96}\\
\toprule Risk & Left minus right & \shortstack{Paired\\midpoint} & \shortstack{Numerical\\half-width} & \shortstack{Simultaneous\\interval}\\
\midrule\endfirsthead
\toprule Risk & Left minus right & \shortstack{Paired\\midpoint} & \shortstack{Numerical\\half-width} & \shortstack{Simultaneous\\interval}\\
\midrule\endhead
0.50 & Online risk-DP minus Quadratic & -0.0001514 & 0.0000011 & $[-0.0003967,0.0000939]$ \\
0.50 & Online risk-DP minus ReLU & -0.0001971 & 0.0000011 & $[-0.0004424,0.0000481]$ \\
0.50 & Online risk-DP minus Adaptive spline & -0.0002038 & 0.0000011 & $[-0.0004491,0.0000414]$ \\
0.50 & Quadratic minus ReLU & -0.0000457 & 0.0000011 & $[-0.0002910,0.0001995]$ \\
0.50 & Quadratic minus Adaptive spline & -0.0000524 & 0.0000011 & $[-0.0002977,0.0001928]$ \\
0.50 & ReLU minus Adaptive spline & -0.0000067 & 0.0000011 & $[-0.0002520,0.0002386]$ \\
0.75 & Online risk-DP minus Quadratic & -0.0001538 & 0.0000012 & $[-0.0003991,0.0000915]$ \\
0.75 & Online risk-DP minus ReLU & -0.0001628 & 0.0000012 & $[-0.0004081,0.0000825]$ \\
0.75 & Online risk-DP minus Adaptive spline & -0.0002433 & 0.0000012 & $[-0.0004887,0.0000020]$ \\
0.75 & Quadratic minus ReLU & -0.0000090 & 0.0000012 & $[-0.0002543,0.0002363]$ \\
0.75 & Quadratic minus Adaptive spline & -0.0000896 & 0.0000012 & $[-0.0003349,0.0001557]$ \\
0.75 & ReLU minus Adaptive spline & -0.0000805 & 0.0000012 & $[-0.0003258,0.0001648]$ \\
\bottomrule\end{longtable}\endgroup
\noindent The sampling radius is the predeclared $1/4096$ for every contrast; the family has coverage at least 95 percent. The exact comparison margin is $1/1280$, and the code rounds it inward when classifying a contrast.

The retained execution contains 32,768 arm--state records. Of the twelve fixed contrasts, 12 establish equivalence at the declared margin and 0 exclude zero. The largest simultaneous half-width is 0.00024531. There are 4750 rejected first proposals and 0 entries to the final exact fallback. These counts refer to the new complete services, not the historical R95 actors.
}{\paragraph{Execution record.}The source-first protocol and executable analysis are committed before confirmation. No numerical outcome is asserted in this source snapshot; a completed execution must supply the current source-bound tables.}

\paragraph{Construction and complete work.}
Label construction, all feature-search fits, state-cover qualification, calibration, serialization and reload are recorded separately for each representation. A model's standalone purchase bears the full shared label cost rather than an arbitrary one-third allocation. Five fresh processes per arm and risk level evaluate the same fixed 32-state timing workload; the parent timer includes interpreter startup, identity checking, loading, qualification, output and shutdown. CPU frequency is recorded, not controlled. Timing ranges are descriptive and feed a conditional sensitivity region, not a confidence interval for future hardware. One-time independent audit cost is shown separately and can be included explicitly in the acquisition charge.

\paragraph{Prospective certificate ablation.}
The builders of the global norm, signed scalar and joint local-support bounds are fixed before confirmation, with the same residual cover and own-continuation resolution. Their proposed-policy identities remain explicit. A scalar bound for a rejected fitted proposal is not attached to the service's different fallback action. The comparison records qualification failures and paid construction work, while distinguishing these fixed builders from a separately tuned solver frontier. The earlier screening experiment and its negative net-time result remain unchanged.

\paragraph{Independent endpoints.}
The independent implementation uses exact rational arithmetic for all paired means and coverage constants. It does not import the production statistical formula. It also chooses eight random rows from each arm--risk stratum, reconstructs the stored candidate actions, and forms a finer original-law complete-policy enclosure and an original-action lower witness from separately written economic primitives. The finer interval must be contained in the reported one. Conventional proposal selection and the full fitting stack are not independently replayed. Sixty-four random endpoint checks are more informative than a fixed hand-selected fixture catalogue, but they are not an independent calculation of all 32,768 endpoints.

The experiment retains its precise state, control and horizon scope. The product quadrature result applies analytically to two innovation coordinates under its convexity hypotheses. It is not a completed two-control, two-innovation solver tournament, and the inherited adverse vector frontier remains visible. The revised comparison addresses the precision and identification questions within the original NBO program without borrowing evidence from an unexecuted model.

\section{Acquiring and qualifying neural readouts}\label{sec:study97}
The new experiment asks whether a construction directed at a valid residual account improves the delivered service, and whether any improvement is attributable to neural geometry. The two questions require separate controls. There are seven services: conventional online risk-DP, and least-squares (LS) and certificate-oriented (CA) readouts for each of the quadratic, oblique-ReLU, and coordinate-spline families. The risk levels remain $1/2$ and $3/4$. Design ranks are four, five, and six; training-label budgets are 256 and 1,024. The resulting 84 configurations are fixed in a remotely committed protocol before the campaign.

Within a family--budget cell the two readouts use identical selected geometry and have the same available original-model labels. Both pay for the 1,089-site cover, including repeated primitive queries used by whole-domain validation. The CA procedure solves \eqref{eq:readout97}, rounds its actual coefficients, and applies Proposition~\ref{prop:readout97}. The family tag, feature directions, offsets, and readout coefficients occupy 64, 96, or 128 bytes at ranks four, five, or six, respectively. All feature-search trials and rank rejections remain in the receipts; matching design rank does not imply equal conditioning or equal fitting work.

Each cell has 16 service-validation states and 64 paired diagnostic states. The complete controller and admission targets are unchanged from R96, not retuned to make CA appear successful. Its stage record distinguishes proposal generation, the first original-law certificate, repair generation, the repair certificate, quadrature refinement, and rational fallback. In particular, the quality of an accepted repaired service is not attributed to its rejected neural proposal.

Across the 84 fixed configurations, all 5,376 complete services meet the unchanged admission criteria. There are 1,401 rejected first proposals, 0 refined rows, and 0 exact-fallback entries. The verified readout rule strictly reduces the whole-domain residual span in 36 of the 36 matched family--budget comparisons. The predeclared price grid contains 1,800 comparisons: 156 conditional adoption regions, 199 conditional rejection regions, and 1,445 unresolved regions. These counts are sensitivity classifications, not frequencies of economically calibrated opportunities.

\begin{table}[t]
\centering\small
\caption{Complete-service diagnostics over the fixed rank and label budgets}\label{tab:diagnostics97}
\begin{tabular}{llrrrr}\toprule
Risk & Service & Proposals & Rejected & Refined & Exact \\ \midrule
0.5 & Online & 384 & 0 & 0 & 0 \\
0.5 & Quadratic LS & 384 & 39 & 0 & 0 \\
0.5 & Quadratic CA & 384 & 101 & 0 & 0 \\
0.5 & ReLU LS & 384 & 127 & 0 & 0 \\
0.5 & ReLU CA & 384 & 135 & 0 & 0 \\
0.5 & Spline LS & 384 & 135 & 0 & 0 \\
0.5 & Spline CA & 384 & 116 & 0 & 0 \\
0.75 & Online & 384 & 0 & 0 & 0 \\
0.75 & Quadratic LS & 384 & 52 & 0 & 0 \\
0.75 & Quadratic CA & 384 & 109 & 0 & 0 \\
0.75 & ReLU LS & 384 & 159 & 0 & 0 \\
0.75 & ReLU CA & 384 & 152 & 0 & 0 \\
0.75 & Spline LS & 384 & 132 & 0 & 0 \\
0.75 & Spline CA & 384 & 144 & 0 & 0 \\
\bottomrule\end{tabular}
\par\smallskip\parbox{\linewidth}{\footnotesize Each risk--service row aggregates three design ranks and two training budgets, with 64 states per cell. Rejected first proposals remain in the complete service. These are construction diagnostics, not a new population superiority test.}
\end{table}


\paragraph{Residual improvement and delivered decisions.}
Table~\ref{tab:budgets97} retains every matched cell, rather than selecting a favorable rank or risk level. The important comparison is between the certified span and the rejected-proposal count \emph{within} that cell. A smaller global span is an improvement in the attained residual account; it need not improve the action-relevant discrepancy on the visited states. Thus a rejected CA proposal is evidence about the alignment of this objective with the delivered decision, not evidence that its smaller validated span was fictitious. The complete LS and CA outcomes, including reversals, remain visible. No confidence interval from R96 is transferred to these newly fitted actors.

\paragraph{Acquisition and a purchase.}
A fresh acquisition process starts from economic primitives. Its internal construction timer supplies a lower endpoint for the measured process cost; the full process, including imports, file output, durable serialization and a subsequent diagnostic probe, supplies a conservative upper endpoint. This deliberately broad bracket avoids subtracting a noisy probe timer to manufacture a small setup cost. It includes repeated labels and any unsuccessful candidate construction. Research authoring and the separate referee reimplementation are not silently part of a purported deployment microbenchmark: their execution receipts are provided separately, and a purchaser requiring an additional audit must add its charge to acquisition or the fixed fee.

For rank six and 1,024 labels, three fresh processes per arm and risk level execute the same 32-state workload. Table~\ref{tab:clocks97} reports all clock ranges. The price grid fixes $m\in\{1,10,100,1000,10000\}$, $\lambda\in\{0,0.01,0.1,1,10\}$, and $F\in\{0,0.01,0.1\}$. Each candidate is compared with online risk-DP and with quadratic LS. We use $[d^-,d^+]=[-1/256,1/256]$ for the new complete services, rather than estimating a population sign from 64 diagnostic states. Acquisition intervals are subtracted as $[K_a^- -K_b^+,K_a^+-K_b^-]$. Equations~\eqref{eq:gainlower97}--\eqref{eq:gainupper97} therefore classify both favorable and unfavorable regions without dropping the factor of 32. These are conditional deployment comparisons in normalized model units, not a calibrated welfare distribution.
\begin{table}[t]
\centering\small
\caption{Measured acquisition brackets and whole-workload clocks}\label{tab:clocks97}
\begin{tabular}{llrrr}\toprule
Risk & Service & Acquisition bracket & Workload median & Workload range \\ \midrule
0.5 & Online & [0.185, 1.306] & 0.841 & [0.826, 0.843] \\
0.5 & Quadratic LS & [0.110, 0.755] & 0.589 & [0.588, 0.606] \\
0.5 & Quadratic CA & [0.216, 1.011] & 0.694 & [0.675, 0.712] \\
0.5 & ReLU LS & [0.218, 1.060] & 0.758 & [0.754, 0.766] \\
0.5 & ReLU CA & [0.236, 1.057] & 0.749 & [0.735, 0.756] \\
0.5 & Spline LS & [0.191, 1.035] & 0.768 & [0.761, 0.773] \\
0.5 & Spline CA & [0.225, 1.084] & 0.753 & [0.744, 0.754] \\
0.75 & Online & [0.201, 1.394] & 0.849 & [0.838, 0.868] \\
0.75 & Quadratic LS & [0.114, 0.815] & 0.580 & [0.564, 0.606] \\
0.75 & Quadratic CA & [0.231, 1.013] & 0.678 & [0.676, 0.696] \\
0.75 & ReLU LS & [0.189, 1.068] & 0.759 & [0.707, 0.811] \\
0.75 & ReLU CA & [0.224, 1.101] & 0.731 & [0.702, 0.732] \\
0.75 & Spline LS & [0.197, 1.099] & 0.734 & [0.717, 0.775] \\
0.75 & Spline CA & [0.231, 1.163] & 0.729 & [0.726, 0.746] \\
\bottomrule\end{tabular}
\par\smallskip\parbox{\linewidth}{\footnotesize Seconds. Rank six and 1,024 training labels. Acquisition lower endpoint is the timed internal construction; the upper endpoint is its entire fresh process including the subsequent diagnostic probe. Each workload has 32 decisions. Three fresh workload repetitions are descriptive, with frequency uncontrolled. Neither bracket is a statistical confidence interval for future execution.}
\end{table}


\paragraph{Reconstruction and independent verification.}
The original R96 ReLU and spline payloads were reconstructed byte for byte in the development check, whereas quadratic feature selection produced different bytes. The check is retained, not overwritten by the new tie convention. An exact expansion of each quadratic into monomial coefficients supplies a uniform function-difference bound on the unit square. A small function difference does not establish identical actions or authorize replacing the frozen R96 model in its confirmation catalogue.

The new independent proposal implementation derives the continuous terminal-risk polynomial from economic primitive evaluations at three innovation points. It implements its own tail integration, envelope derivative, stable discrete tail selection, and action bisection; it imports none of the production interval, actor, quadrature, or fitting modules. Its binary64 proposal comparison is distinct from the exact-rational endpoint and inference audit. The audit receipt records every discrepancy and the explicit action-quantization tolerance. The independently reimplemented acquisition optimizer remains a separate verification obligation; fitting-stack independence is not inferred from an independent proposal calculation.

The previously separate ARM64 evidence is retained under a single pinned lineage. It replays the original frozen R96 models and the same initial states, reports identical production CSV bytes, and adds zero independent economic observations. It is not a refit of the new models, a full ARM64 acquisition benchmark, or an energy measurement. The new construction campaign and this historical architecture replay are not pooled.

The state, action, innovation, and horizon scope of the continuous investment tournament remains explicit. The product quadrature theorem and the retained two-control analysis do not become a completed longer-horizon vector-risk tournament by changing their names. The new evidence addresses readout acquisition, service work, and purchase resolution within the original NBO model; the complete inherited vector and horizon results remain available under their own assumptions.


\section{Conclusion}
Neural Bellman Operators construct attained continuations and feasible dynamic policies with an original-optimum error account. Original-owner compilation and reference-centered recursive-risk qualification preserve the economic object while exposing the work required to evaluate it. The continuous investment construction supplies the necessary curvature, supports and full-law inequalities from economic primitives.

The additional action-aligned result identifies the quantities that control a proposal's decision gap: the width of an original-objective enclosure, the directional support mismatch and the constrained stationarity residual. Prospective selection can then target rejection or complete work rather than infer either from a smaller global residual span. The exact purchase mechanism demonstrates a strictly positive economic value of an oblique continuation relative to the best specified additive and quadratic classes, while equal conventional geometry attains equal value. Nonuniform-law and explicitly quantified randomized information results place the precision account on a wider mathematical basis.

The investment tournament remains an essential discipline on these constructions. Its quadratic least-squares winner, increased aggregate rejections under residual-oriented fitting, and absence of certified neural adoption over that winner are preserved. The mathematical mechanisms and that empirical comparison answer different questions. Together they identify what an NBO must construct, what its certificate must cover and what evidence is needed to establish a useful decision at a competitive complete cost.

\bibliographystyle{ecta-fullname}
\bibliography{references99}
\end{document}
